\section*{Bab \ref{ch:b3:membrane-traffic} --- Lalu Lintas Membran dan Penyortiran Protein}

\begin{solution}{exo:b3:membrane-traffic:1}
Bentangan amino-terminal yang hidrofob: \omterm{def:b3:membrane-traffic:signals}{peptida isyarat}, menuju ke dalam
retikulum endoplasma (lalu secara baku terus ke permukaannya). PKKKRKV:
\omterm{def:b3:membrane-traffic:signals}{isyarat penempatan nukleus}, menuju ke dalam nukleus lewat porinya.
Heliks amfipatik amino-terminal: \omterm{def:b3:membrane-traffic:signals}{praurutan} mitokondria, menuju ke dalam
matriks mitokondria. KDEL: pengambilan kembali dari Golgi ke
retikulumnya. Manosa-6-fosfat: dari trans-Golgi menuju \omterm{def:b3:membrane-traffic:endocytosis}{endosom}
dan \omterm{def:b3:membrane-traffic:endocytosis}{lisosom}.
\end{solution}

\begin{solution}{exo:b3:membrane-traffic:2}
(1) Tanpa membran, rantai yang ditranslasikan lebih panjang daripada protein
sekresinya kira-kira dua puluh residu: sebuah segmen dibuang selama
sekresinya. (2) Dengan membran yang hadir selama translasinya, hasilnya
berukuran matang dan terlindung dari protease: hasilnya telah masuk ke
vesikelnya lalu kehilangan segmennya di dalamnya. (3) Membran yang ditambahkan
sesudah translasinya tidak berbuat keduanya: masuknya terkopel dengan
sintesisnya --- kotranslasi --- dan segmen itu, yaitu \omterm{def:b3:membrane-traffic:signals}{peptida isyarat}, hanya
bekerja pada rantai yang nasen.
\end{solution}

\begin{solution}{exo:b3:membrane-traffic:3}
Retikulum ke Golgi: \omterm{def:b3:membrane-traffic:coats}{COPII}, maju. Golgi ke retikulum: \omterm{def:b3:membrane-traffic:coats}{COPI},
mundur. Trans-Golgi ke \omterm{def:b3:membrane-traffic:endocytosis}{endosom}: \omterm{def:b3:membrane-traffic:coats}{klatrin} beserta adaptornya,
maju menuju \omterm{def:b3:membrane-traffic:endocytosis}{lisosom}. Membran plasma ke \omterm{def:b3:membrane-traffic:endocytosis}{endosom}: \omterm{def:b3:membrane-traffic:coats}{klatrin},
ke dalam (endositosis).
\end{solution}

\begin{solution}{exo:b3:membrane-traffic:4}
Gula hanya ditambahkan di dalam lumen retikulum dan Golginya. Muka
lumenal setiap vesikel tetap lumenal sepanjang pertunasan dan fusinya,
dan ketika sebuah vesikel berfusi dengan membran plasmanya, muka lumenalnya
menjadi bagian luar selnya; orientasinya tak pernah terbalik, sehingga
apa yang terglikosilasi di dalam berakhir di luar.
\end{solution}

\begin{solution}{exo:b3:membrane-traffic:5}
$t^{*} = \ln(0.05/0.2)/(0.05 - 0.2) = \ln 4/0.15 = \qty{9.2}{min}$;
$B(t^{*}) = (0.2/(-0.15))(e^{-1.85} - e^{-0.46}) = 0.63$. Pada \qty{60}{min}
$A \approx 0$, $B = 1.33\,(e^{-3} - e^{-12}) = 0.066$, sehingga $C = 0.93$.
\end{solution}

\begin{solution}{exo:b3:membrane-traffic:6}
\omterm{def:b3:membrane-traffic:signals}{Peptida isyaratnya} menaruh ujung aminonya di dalam lumennya; tiap
bentangan hidrofobnya secara berselang-seling menjadi henti-pindah atau
mulai-pindah, sehingga ada empat heliks lintas membran (kira-kira residu 30--52,
90--112, 150--172, 210--232). Segmennya: ujung aminonya lumenal;
52--90 sitosolik; 112--150 lumenal; 172--210 sitosolik; ujung
karboksinya lumenal. Glikosilasinya hanya pada segmen lumenal --- yaitu
ujung aminonya, 112--150, dan ujung karboksinya --- yang menjadi
ekstrasel pada permukaannya.
\end{solution}

\begin{solution}{exo:b3:membrane-traffic:7}
$J_{\max} = \num{30000}\times 0.25\times 0.05/0.30 = 1250$ tiap menit;
bagian permukaannya $k_{r}/(k_{i} + k_{r}) = 0.05/0.30 = 0.17$. Melipatduakan $R$
melipatduakan $J$ secara persis; laju kembalinya adalah lehernya (sebagian
besar reseptornya berada di dalam), dan melipatduakan $k_{r}$ akan memberi $\num{30000}\times 0.25
\times 0.1/0.35 \approx 2140$, yakni faktor $1.7$.
\end{solution}

\begin{solution}{exo:b3:membrane-traffic:8}
(a) Tanpa reseptor: hidrolase yang bertanda itu mengikuti rute bakunya lalu
disekresikan; \omterm{def:b3:membrane-traffic:endocytosis}{lisosomnya} kekurangan enzim lalu terisi bahan yang tak
tercerna --- sebuah penyakit penimbunan. (b) Tanpa fosfotransferase: fenotipe
yang sama lewat cacat yang berbeda (penyakit sel-I), dengan enzimnya
tak bertanda lalu disekresikan. (c) \omterm{def:b3:membrane-traffic:endocytosis}{Endosom} yang dinetralkan: reseptornya tak
dapat melepaskan muatannya, sehingga ia tidak didaur ulang dan hidrolasenya
salah antar atau disekresikan; LDL dan reseptor lain pun berhenti berdaur.
\end{solution}

\begin{solution}{exo:b3:membrane-traffic:9}
Tirosinanya termasuk \omterm{def:b3:bioinformatics:motif}{motif} ekor yang dikenali adaptor \omterm{def:b3:membrane-traffic:coats}{klatrin};
tanpa tirosina itu reseptornya tidak terkumpul ke dalam lubang bersalut.
Reseptornya mengikat LDL secara normal tetapi tersebar merata di seluruh
permukaannya, tersingkir dari lubangnya, dan ligannya tinggal di luar --- sebuah
penyakit penyortiran, bukan penyakit pengikatan.
\end{solution}

\begin{solution}{exo:b3:membrane-traffic:10}
Dengan $k_{1} = k_{2} = k$, cobalah $B = kt\,e^{-kt}$: $\mathrm{d}B/\mathrm{d}t
= k e^{-kt} - k^{2}t e^{-kt} = kA - kB$, seperti yang dituntut, dengan $B(0) = 0$.
Puncaknya ketika $1 - kt = 0$: $t^{*} = 1/k$, $B = e^{-1} = 0.37$. Secara umum
$C(t) = 1 - (k_{2}e^{-k_{1}t} - k_{1}e^{-k_{2}t})/(k_{2} - k_{1})$,
yang setangkup di bawah pertukaran $k_{1}$ dan $k_{2}$: kurva butirannya
saja tak dapat mengatakan kompartemen yang mana yang cepat.
\end{solution}

\begin{solution}{exo:b3:membrane-traffic:11}
Tiap vesikelnya berluas $4\pi(50\,\text{nm})^{2} = \qty{0.031}{\micro
m^2}$; seribu tiap menit berarti \qty{31}{\micro m^2/min}, sehingga seluruh
permukaannya dalam $3000/31 \approx \qty{95}{min}$. Selnya tidak mengerut
karena membrannya kembali lewat eksositosis (vesikel daur ulang) dengan laju
yang sama --- yakni keseimbangan aliran. Tiap vesikelnya memuat $5\times 10^{-19}$ L;
seribu tiap menit selama sejam berarti $3\times 10^{-14}$ L, kira-kira
\qty{1.5}{\%} volume sel \qty{2}{pL} tiap jam.
\end{solution}

\begin{solution}{exo:b3:membrane-traffic:12}
Toksin botulinum bekerja di dalam terminal motornya: tak ada asetilkolina yang
dilepaskan, ototnya tak menerima perintah lalu terkulai lemas --- kelumpuhan
lemas. Toksin tetanus dibawa naik sepanjang akson motornya lalu menyeberang ke
antarneuron penghambat yang biasanya mengekang neuron motornya; dengan
\omterm{def:b3:membrane-traffic:fusion}{SNARE-nya} terpotong, tak ada glisina atau GABA yang dilepaskan, neuron motornya
menyala tanpa lawan dan setiap ototnya berkerut --- kelumpuhan kaku. Beberapa
nanogram toksin botulinum yang disuntikkan ke dalam otot yang terlalu aktif
membungkamnya secara setempat selama berbulan-bulan (sampai \omterm{def:b3:membrane-traffic:fusion}{SNARE} baru dibuat
dan terminalnya pulih): sebuah pengobatan bagi distonia, kekakuan, juling, dan kerutan.
\end{solution}

\begin{solution}{pb:b3:membrane-traffic:1}
\textbf{1.} $10^{7}\times \qty{25000}{Da}\times 1.66\times 10^{-24}$ g
$= \qty{0.42}{pg}$ tiap menit, \qty{600}{pg} tiap hari --- dua kali kandungan
proteinnya sendiri.
\textbf{2.} Volume vesikelnya $\tfrac{4}{3}\pi\,35^{3} = 1.8\times 10^{5}$
nm$^{3}$, dan \qty{30}{\%} di antaranya $5.4\times 10^{4}$; satu molekul enzim
$\tfrac{4}{3}\pi\,2^{3} = 34$ nm$^{3}$: sekitar \num{1600} molekul tiap
vesikel; $10^{7}/1600 \approx 6200$ vesikel tiap menit.
\textbf{3.} Luasnya $4\pi\,35^{2} = \qty{0.0154}{\micro m^2}$ masing-masing, \qty{95}{\micro
m^2} tiap menit: seluas \qty{6000}{\micro m^2} retikulumnya dalam kira-kira
sejam andaikan tak ada yang kembali.
\textbf{4.} $\qty{95}{\micro m^2} = 9.5\times 10^{7}$ nm$^{2}$, dua lembaran
pada \qty{0.7}{nm^2}: $2.7\times 10^{8}$ lipid tiap menit. Selnya
mengembalikan membrannya lewat vesikel \omterm{def:b3:membrane-traffic:coats}{COPI} dan daur ulang, lalu mensintesis
lipid untuk mengganti yang pergi untuk selamanya di dalam butirannya.
\textbf{5.} Volume butirannya $5.2\times 10^{8}$ nm$^{3}$, dan \qty{30}{\%} di antaranya
$1.6\times 10^{8}$: $4.7\times 10^{6}$ molekul tiap butiran; $6\times
10^{8}$ molekul tiap jam mengisi sekitar $130$ butiran.
\textbf{6.} Tiap butirannya menambahkan $\pi d^{2} = \qty{3.1}{\micro m^2}$; $300$
butiran menambahkan \qty{940}{\micro m^2} pada \qty{50}{\micro m^2}: dua puluh
kali lipat. Membrannya harus diambil kembali lewat endositosis penggenap dalam
hitungan menit.
\textbf{7.} $t^{*} = \ln 2/0.05 = \qty{13.9}{min}$, $B = 0.50$.
\textbf{8.} $A = e^{-3} = 0.05$; $B = 2(e^{-1.5} - e^{-3}) = 0.35$;
$C = 0.60$.
\textbf{9.} \qty{10}{min} dan \qty{20}{min}; totalnya rata-rata \qty{30}{min}.
\textbf{10.} $t^{*} = \ln(0.1)/(-0.09) = \qty{25.6}{min}$; $B(t^{*}) =
(0.1/(-0.09))(e^{-2.56} - e^{-0.256}) = 0.77$: Golginya membengkak,
sarat dengan tiga perempat muatannya.
\textbf{11.} Pada puncaknya $k_{1}e^{-k_{1}t^{*}} = k_{2}e^{-k_{2}t^{*}}$;
dengan mensubstitusikannya ke dalam $B$ diperoleh $B_{\max} = (k_{1}/k_{2})^{\,k_{2}/(k_{2} -
k_{1})}$: pangkatnya adalah $k_{2}/(k_{2} - k_{1})$. Bagi $k_{1} > k_{2}$
pangkatnya negatif dan basisnya di atas $1$; bagi $k_{1} < k_{2}$
basisnya di bawah $1$ dan pangkatnya positif --- pada kedua kasus nilainya
berada di bawah $1$ (periksalah: $2^{-1} = 0.5$ di sini).
\textbf{12.} Bagi $k_{1} \gg k_{2}$, $B(t) \to e^{-k_{2}t}$: labelnya
langsung berada di dalam Golginya lalu pergi pada $k_{2}$ --- langkah
retikulumnya tak terlihat dan kurva Golginya menjadi peluruhan sederhana.
\textbf{13.} $\theta = 2/(2+2) = 0.5$; $J = \num{50000}\times 0.2\times
0.1\times 0.5/(0.1 + 0.1) = 2500$ butir tiap menit, \num{150000} tiap
jam.
\textbf{14.} $1.5\times 10^{5}\times 1500 = 2.3\times 10^{8}$ molekul
kolesterol tiap jam; $5\times 10^{9}/2.3\times 10^{8} \approx 22$ jam.
\textbf{15.} Bagian permukaannya $k_{r}/(k_{r} + k_{i}\theta) = 0.5$; satu
perjalanan bolak-balik memakan $1/(k_{i}\theta) + 1/k_{r} = 10 + 10 = \qty{20}{min}$:
sekitar $60$ perjalanan dalam dua puluh jam.
\textbf{16.} $J = 1250$ tiap menit. Di bawah kejenuhannya $J \propto R\,[L]$,
sehingga dengan separuh reseptornya LDL plasmanya harus berlipat dua demi
pembersihan yang sama.
\textbf{17.} $R$ kembali ke \num{50000}: pembersihannya pulih, LDL plasmanya
turun kembali menuju normal --- itulah pengaruh statinnya.
\textbf{18.} Tanpa reseptor, LDL-nya hanya dibersihkan lewat rute
tak khas yang lambat, sehingga ia beredar berhari-hari alih-alih berjam-jam,
teroksidasi, lalu diserap reseptor pemulung makrofag,
yang menjadi sel busa plak aterosklerosis; dan statin,
yang bekerja dengan mengimbas reseptor, tak punya apa pun untuk diimbas.
\textbf{19.} $V = \tfrac{4}{3}\pi(0.25)^{3} = \qty{0.065}{\micro m^3} =
6.5\times 10^{-17}$ L. Pada pH $7.2$: $6.3\times 10^{-8}\times 6.5\times
10^{-17} = 4\times 10^{-24}$ mol --- kira-kira $2$ proton bebas. Pada pH $4.7$:
$2\times 10^{-5}\times 6.5\times 10^{-17} = 1.3\times 10^{-21}$ mol,
sekitar $780$.
\textbf{20.} $50\times 780 \approx \num{39000}$ proton; $50$ pompa pada
$200$ tiap detik memindahkan \num{10000} tiap detik: sekitar \qty{4}{s}.
\textbf{21.} Memompa muatan positif ke dalam membuat lumennya positif, dan
beberapa ribu muatan akan menghentikan pompanya; klorida masuk lewat sebuah
saluran (atau kation keluar) untuk menetralkan muatannya.
\textbf{22.} Enzim yang tak aktif pada pH netral tak merugikan bila sebuah
\omterm{def:b3:membrane-traffic:endocytosis}{lisosom} bocor atau bila enzimnya salah sortir ke sitosol atau ke permukaannya:
tuntutan asamnya mengurung pencernaannya di dalam kompartemen yang dibangun
untuk itu.
\textbf{23.} Keduanya pasti mengusung sebuah pengindera pH --- histidina, yang
protonasinya di dekat pH~6 mengubah permukaan pengikatnya: reseptor LDL
melipat sebuah \omterm{def:b3:structural-biology:levels}{domain} menutupi tapak pengikatnya sendiri ketika terprotonasi,
dan reseptor manosa-6-fosfat kehilangan kegemarannya dengan cara serupa.
\textbf{24.} Pemekatan $10^{7.2 - 4.7} = 10^{2.5} \approx 300$ kali lipat;
basa yang terperangkap itu memakan proton, menaikkan pH \omterm{def:b3:membrane-traffic:endocytosis}{lisosomnya}, lalu
menonaktifkan hidrolasenya --- dan begitulah klorokuin membunuh parasit malaria
di dalam vakuola makanannya serta menghalangi \omterm{def:b3:membrane-traffic:endocytosis}{autofagi} pada
percobaan.
\textbf{25.} Puncak Golginya pada \qty{14}{min}; sekitar \num{6200} vesikel
meninggalkan retikulumnya tiap menit; \num{150000} butir LDL terserap tiap
jam.
\end{solution}
